\documentclass{article}
\usepackage{amsmath,amssymb,amsthm}
\usepackage{algorithm}
\usepackage{algpseudocode}
\usepackage{geometry}
\geometry{margin=1in}
\newtheorem{theorem}{Theorem}
\newtheorem{lemma}{Lemma}
\newtheorem{assumption}{Assumption}
\newcommand{\E}{\mathbb{E}}
\newcommand{\R}{\mathbb{R}}
\newcommand{\norm}[1]{\left\|#1\right\|}
\newcommand{\inner}[2]{\left\langle #1,#2\right\rangle}
\newcommand{\argmin}{\operatorname*{arg\,min}}

\begin{document}

\section*{Local SGD (Federated Averaging)}

\section*{Mathematical Formulation}
Consider a collection of $K$ clients.  Each client $k\in\{1,\dots,K\}$ possesses a local convex, $L$‑smooth loss
\[
f_k(w)=\E_{z\sim\mathcal{D}_k}\big[\ell(w;z)\big],\qquad w\in\R^d .
\]
The global objective is the average
\[
f(w)=\frac1K\sum_{k=1}^K f_k(w).
\]

The algorithm proceeds in \emph{communication rounds} $t=0,1,\dots,T-1$.  
At the beginning of round $t$ a \emph{central model} $w_t$ is broadcast to all clients.  
Each client performs $E$ local SGD steps with stepsize $\eta>0$:
\[
\begin{aligned}
w_{t}^{k,0}&=w_t,\\
w_{t}^{k,e+1}&= w_{t}^{k,e}-\eta\, g_{t}^{k,e},
\qquad e=0,\dots,E-1,
\end{aligned}
\tag{1}
\]
where $g_{t}^{k,e}$ is a stochastic gradient of $f_k$ at $w_{t}^{k,e}$,
\[
\E\!\left[g_{t}^{k,e}\mid w_{t}^{k,e}\right]=\nabla f_k\!\left(w_{t}^{k,e}\right).
\]

After the $E$ local steps the client sends its final iterate $w_{t}^{k,E}$ back to the server, which averages:
\[
w_{t+1}= \frac1K\sum_{k=1}^K w_{t}^{k,E}.
\tag{2}
\]

Define the \emph{client drift} at round $t$ as
\[
\delta_t \;:=\; \frac1K\sum_{k=1}^K\big\|\nabla f_k(w_t)-\nabla f(w_t)\big\|.
\tag{3}
\]

\section*{Pseudocode}
\begin{algorithm}[H]
\caption{Local SGD (FedAvg)}
\begin{algorithmic}[1]
\State \textbf{Input:} stepsize $\eta$, local steps $E$, total rounds $T$, initial model $w_0$.
\For{$t=0$ \textbf{to} $T-1$}
    \State Server broadcasts $w_t$ to all $K$ clients.
    \For{each client $k$ \textbf{in parallel}}
        \State $w_{t}^{k,0}\gets w_t$
        \For{$e=0$ \textbf{to} $E-1$}
            \State Sample minibatch $\mathcal{B}_{t}^{k,e}$ from $\mathcal{D}_k$.
            \State $g_{t}^{k,e}\gets \frac1{|\mathcal{B}_{t}^{k,e}|}\sum_{z\in\mathcal{B}_{t}^{k,e}}\nabla\ell(w_{t}^{k,e};z)$.
            \State $w_{t}^{k,e+1}\gets w_{t}^{k,e}-\eta\, g_{t}^{k,e}$.
        \EndFor
        \State Send $w_{t}^{k,E}$ to server.
    \EndFor
    \State Server aggregates: $w_{t+1}\gets \frac1K\sum_{k=1}^K w_{t}^{k,E}$.
\EndFor
\State \textbf{Output:} $\bar w_T:=\frac1T\sum_{t=0}^{T-1} w_t$ (or the last $w_T$).
\end{algorithmic}
\end{algorithm}

\section*{Dry Run Example}
Minimize $f(w)=(w-3)^2$ (convex, $L=2$).  Use a single client ($K=1$), $E=1$, stepsize $\eta=0.1$, and start $w_0=0$.
The stochastic gradient equals the true gradient $g_t=2(w_t-3)$.

\begin{center}
\begin{tabular}{c|c|c|c}
$t$ & $w_t$ & $g_t=2(w_t-3)$ & $f(w_t)$\\ \hline
0 & $0$ & $-6$ & $9$\\
1 & $w_1=0-0.1(-6)=0.6$ & $2(0.6-3)=-4.8$ & $(0.6-3)^2=5.76$\\
2 & $w_2=0.6-0.1(-4.8)=1.08$ & $2(1.08-3)=-3.84$ & $(1.08-3)^2=3.6864$\\
3 & $w_3=1.08-0.1(-3.84)=1.464$ & $2(1.464-3)=-3.072$ & $(1.464-3)^2=2.3677$
\end{tabular}
\end{center}
The objective value decreases at each iteration, illustrating the local update (1) and the trivial averaging (2).

\newpage
\section*{Assumptions}
\begin{assumption}[Convexity and Smoothness]\label{as:convex}
For every client $k$,
\[
f_k \text{ is convex and } L\text{-smooth: } 
f_k(y)\le f_k(x)+\inner{\nabla f_k(x)}{y-x}+\frac{L}{2}\norm{y-x}^2,\;\forall x,y\in\R^d .
\]
\end{assumption}
\begin{assumption}[Unbiased Stochastic Gradients]\label{as:unbiased}
For all $k,e,t$,
\[
\E\!\big[g_{t}^{k,e}\mid w_{t}^{k,e}\big]=\nabla f_k\!\big(w_{t}^{k,e}\big).
\]
\end{assumption}
\begin{assumption}[Bounded Variance]\label{as:variance}
There exists $\sigma^2>0$ such that
\[
\E\!\big[\|g_{t}^{k,e}-\nabla f_k(w_{t}^{k,e})\|^2\mid w_{t}^{k,e}\big]\le\sigma^2,\qquad\forall k,e,t.
\]
\end{assumption}
\begin{assumption}[Bounded Gradient Norm]\label{as:gradbound}
\[
\|\nabla f_k(w)\|\le G,\qquad\forall k,w.
\]
\end{assumption}
\begin{assumption}[Client Drift]\label{as:drift}
There exists $\Delta\ge0$ such that for any $w$,
\[
\frac1K\sum_{k=1}^K\big\|\nabla f_k(w)-\nabla f(w)\big\|^2\le\Delta^2 .
\]
\end{assumption}
All hyper‑parameters satisfy $0<\eta\le \frac{1}{L}$ and $E\ge1$.

\section*{Main Convergence Theorem}
\begin{theorem}\label{thm:conv}
Let the iterates $\{w_t\}$ be generated by Local SGD with the above assumptions.
Define the averaged model $\bar w_T:=\frac1T\sum_{t=0}^{T-1}w_t$.
Then for any optimal solution $w^\star\in\arg\min f$,
\[
\E\!\big[f(\bar w_T)-f(w^\star)\big]\;\le\;
\frac{\norm{w_0-w^\star}^2}{2\eta KET}
\;+\;
\frac{\eta L\sigma^2}{2K}
\;+\;
\frac{\eta L\Delta^2 (E-1)}{2}.
\tag{4}
\]
Consequently, choosing $\eta=\Theta\!\big(\frac{1}{\sqrt{T}}\big)$ yields the rate
\[
\E\!\big[f(\bar w_T)-f(w^\star)\big]=O\!\Big(\frac{1}{\sqrt{T}}+\frac{E-1}{\sqrt{T}}\Big).
\]
\end{theorem}

\section*{Theoretical Properties}
\begin{itemize}
\item \textbf{Stability w.r.t. stepsize.}  The bound (4) holds for any $\eta\le 1/L$, reflecting the standard $L$‑smoothness condition.
\item \textbf{Effect of client drift.}  The term $\frac{\eta L\Delta^2 (E-1)}{2}$ quantifies the penalty incurred by performing multiple local steps; if the data are homogeneous ($\Delta=0$) the penalty disappears and the rate matches centralized SGD.
\item \textbf{Communication‑efficiency.}  Increasing $E$ reduces the number of communication rounds by a factor $E$ while only adding a linear drift penalty, explaining the empirical benefit of FedAvg when $\Delta$ is small.
\item \textbf{Comparison to baselines.}
    \begin{enumerate}
        \item \emph{Centralized SGD} ($K=1$, $E=1$) reduces (4) to the classical $O(1/(\eta T))+\eta\sigma^2$ bound.
        \item \emph{Parallel SGD} ($E=1$, $K>1$) improves the variance term by a factor $K$ (the $\sigma^2/K$ term), matching standard mini‑batch analysis.
    \end{enumerate}
\end{itemize}

\section*{Discussion}
The theorem demonstrates that Local SGD inherits the $O(1/T)$ convergence of centralized SGD up to an additive drift term that scales with $E-1$ and the heterogeneity constant $\Delta$.  In practice the heterogeneity is often modest, so large $E$ yields substantial communication savings with only a mild slowdown.  The proof below is fully constructive and annotates each non‑trivial manipulation with a step number for easy verification.

\newpage
\section*{Preliminaries (Definitions and Lemmas)}
\begin{lemma}[Smoothness inequality]\label{lem:smooth}
If $f$ is $L$‑smooth then for any $x,y\in\R^d$,
\[
f(y)\le f(x)+\inner{\nabla f(x)}{y-x}+\frac{L}{2}\norm{y-x}^2 .
\tag{L1}
\]
\end{lemma}
\begin{lemma}[Convexity inequality]\label{lem:convex}
If $f$ is convex then for any $x,y$,
\[
f(x)-f(y)\le\inner{\nabla f(x)}{x-y}.
\tag{L2}
\]
\end{lemma}
\begin{lemma}[Bound on drift accumulation]\label{lem:drift}
Under Assumption~\ref{as:drift},
\[
\frac1K\sum_{k=1}^K\big\|\nabla f_k(w_{t}^{k,e})-\nabla f(w_t)\big\|^2
\le \Delta^2 + L^2\sum_{e'=0}^{e-1}\norm{w_{t}^{k,e'}-w_t}^2 .
\tag{L3}
\]
\end{lemma}
\begin{proof}
Apply triangle inequality and $L$‑smoothness to each term $\nabla f_k(w_{t}^{k,e})-\nabla f_k(w_t)$, then square and use Assumption~\ref{as:drift}.  Details are omitted for brevity because the inequality will be used only as an upper bound. \hfill $\square$
\end{proof}

\section*{Main Proof (Step‑by‑Step Derivation)}
We prove Theorem~\ref{thm:conv}.  Throughout the proof $\E_t[\cdot]$ denotes expectation conditioned on all randomness up to the beginning of round $t$ (i.e., on $w_t$).

\subsection*{Step 0: Decompose one communication round}
For a fixed round $t$, write the server update (2) as
\[
w_{t+1}=w_t-\eta\underbrace{\frac1{K}\sum_{k=1}^K\frac1E\sum_{e=0}^{E-1} g_{t}^{k,e}}_{:=\; \widetilde g_t}.
\tag{5}
\]

\subsection*{Step 1: Expand squared distance to optimum}
\[
\begin{aligned}
\norm{w_{t+1}-w^\star}^2
&= \norm{w_t-w^\star}^2
-2\eta\inner{\widetilde g_t}{w_t-w^\star}
+\eta^2\norm{\widetilde g_t}^2 .
\end{aligned}
\tag{6}
\]
Taking expectation $\E_t[\cdot]$ and using linearity:
\[
\E_t\norm{w_{t+1}-w^\star}^2
= \norm{w_t-w^\star}^2
-2\eta\,\E_t\!\big[\inner{\widetilde g_t}{w_t-w^\star}\big]
+\eta^2\,\E_t\norm{\widetilde g_t}^2 .
\tag{7}
\]

\subsection*{Step 2: Relate $\widetilde g_t$ to true gradient}
By unbiasedness (Assumption~\ref{as:unbiased}) and the definition of $\widetilde g_t$,
\[
\E_t[\widetilde g_t]=\frac1{K}\sum_{k=1}^K\frac1E\sum_{e=0}^{E-1}
\nabla f_k\!\big(w_{t}^{k,e}\big).
\tag{8}
\]
Add and subtract $\nabla f(w_t)$:
\[
\E_t[\widetilde g_t]=\nabla f(w_t)+\underbrace{\frac1{K}\sum_{k=1}^K\frac1E\sum_{e=0}^{E-1}
\big(\nabla f_k(w_{t}^{k,e})-\nabla f_k(w_t)\big)}_{:=\; \Delta_t}.
\tag{9}
\]

\subsection*{Step 3: Lower bound the inner product term}
Using (9) in the second term of (7) gives
\[
\begin{aligned}
-2\eta\,\E_t\!\big[\inner{\widetilde g_t}{w_t-w^\star}\big]
&= -2\eta\inner{\nabla f(w_t)}{w_t-w^\star}
-2\eta\inner{\Delta_t}{w_t-w^\star}.
\end{aligned}
\tag{10}
\]

Apply Lemma~\ref{lem:convex} (L2) to the first inner product:
\[
\inner{\nabla f(w_t)}{w_t-w^\star}\ge f(w_t)-f(w^\star).
\tag{11}
\]
Thus,
\[
-2\eta\inner{\nabla f(w_t)}{w_t-w^\star}
\le -2\eta\big(f(w_t)-f(w^\star)\big).
\tag{12}
\]

\subsection*{Step 4: Bound the drift inner product}
For the second inner product we use Cauchy–Schwarz and Lemma~\ref{lem:drift}:
\[
\begin{aligned}
\big|\inner{\Delta_t}{w_t-w^\star}\big|
&\le \norm{\Delta_t}\,\norm{w_t-w^\star}\\
&\le \norm{w_t-w^\star}\;
\underbrace{\frac1{K}\sum_{k=1}^K\frac1E\sum_{e=0}^{E-1}
\norm{\nabla f_k(w_{t}^{k,e})-\nabla f_k(w_t)}}_{\le \Delta + L\cdot\frac1E\sum_{e}\norm{w_{t}^{k,e}-w_t}}
\end{aligned}
\tag{13}
\]
Using $ab\le \frac{a^2}{2\beta}+ \frac{\beta b^2}{2}$ with $\beta=1$,
\[
-2\eta\inner{\Delta_t}{w_t-w^\star}
\le \eta\norm{w_t-w^\star}^2 + \eta\Big(\Delta^2+L^2\frac1{E}\sum_{e=0}^{E-1}\norm{w_{t}^{k,e}-w_t}^2\Big).
\tag{14}
\]

\subsection*{Step 5: Upper bound the variance term $\E_t\norm{\widetilde g_t}^2$}
Decompose
\[
\widetilde g_t = \underbrace{\E_t[\widetilde g_t]}_{\text{mean}} + \underbrace{\big(\widetilde g_t-\E_t[\widetilde g_t]\big)}_{\text{zero‑mean noise}} .
\]
Hence,
\[
\E_t\norm{\widetilde g_t}^2
= \norm{\E_t[\widetilde g_t]}^2
+ \E_t\norm{\widetilde g_t-\E_t[\widetilde g_t]}^2 .
\tag{15}
\]

The noise variance is bounded by the per‑client variance (Assumption~\ref{as:variance}):
\[
\E_t\norm{\widetilde g_t-\E_t[\widetilde g_t]}^2
\le \frac1{K^2}\sum_{k=1}^K\frac1{E^2}\sum_{e=0}^{E-1}\E_t\norm{g_{t}^{k,e}-\nabla f_k(w_{t}^{k,e})}^2
\le \frac{\sigma^2}{K E}.
\tag{16}
\]

For the squared mean term we simply drop it (it is non‑negative) to obtain a convenient upper bound:
\[
\norm{\E_t[\widetilde g_t]}^2\le \Big\|\nabla f(w_t)\Big\|^2 + 2\big\|\Delta_t\big\|^2
\le G^2 + 2\Delta^2 + 2L^2\frac1{E}\sum_{e=0}^{E-1}\norm{w_{t}^{k,e}-w_t}^2 .
\tag{17}
\]

Combining (16) and (17) yields
\[
\E_t\norm{\widetilde g_t}^2
\le \frac{\sigma^2}{K E}+ G^2 + 2\Delta^2 + 2L^2\frac1{E}\sum_{e=0}^{E-1}\norm{w_{t}^{k,e}-w_t}^2 .
\tag{18}
\]

\subsection*{Step 6: Relate local iterates to the server model}
From the local update (1) we have for any $e\ge0$,
\[
w_{t}^{k,e+1}-w_t = -\eta\sum_{s=0}^{e} g_{t}^{k,s}.
\tag{19}
\]
Taking norms and using bounded gradients $\|g_{t}^{k,s}\|\le G+\sigma$ (from Assumptions~\ref{as:gradbound} and~\ref{as:variance}) gives
\[
\norm{w_{t}^{k,e}-w_t}\le \eta e (G+\sigma).
\tag{20}
\]
Consequently,
\[
\frac1{E}\sum_{e=0}^{E-1}\norm{w_{t}^{k,e}-w_t}^2
\le \eta^2 (G+\sigma)^2 \frac1{E}\sum_{e=0}^{E-1} e^2
\le \eta^2 (G+\sigma)^2 \frac{(E-1)E(2E-1)}{6E}
\le \eta^2 (G+\sigma)^2 \frac{E^2}{3}.
\tag{21}
\]

\subsection*{Step 7: Assemble the recursion}
Insert the bounds (12)–(14) and (18) into (7):
\[
\begin{aligned}
\E_t\norm{w_{t+1}-w^\star}^2
&\le \norm{w_t-w^\star}^2
-2\eta\big(f(w_t)-f(w^\star)\big) \\
&\quad + \eta\norm{w_t-w^\star}^2
+ \eta\Delta^2
+ \eta L^2 \frac1{E}\sum_{e=0}^{E-1}\norm{w_{t}^{k,e}-w_t}^2 \\
&\quad + \eta^2\Big(\frac{\sigma^2}{K E}+ G^2 + 2\Delta^2 + 2L^2 \frac1{E}\sum_{e=0}^{E-1}\norm{w_{t}^{k,e}-w_t}^2\Big).
\end{aligned}
\tag{22}
\]

Collect the terms that contain $\frac1{E}\sum_{e}\norm{w_{t}^{k,e}-w_t}^2$ and use (21):
\[
\begin{aligned}
\E_t\norm{w_{t+1}-w^\star}^2
&\le \big(1+\eta\big)\norm{w_t-w^\star}^2
-2\eta\big(f(w_t)-f(w^\star)\big) \\
&\quad + \eta\Delta^2
+ \eta^2\Big(\frac{\sigma^2}{K E}+ G^2 + 2\Delta^2\Big) \\
&\quad + \Big(\eta L^2 +2\eta^2 L^2\Big)\eta^2 (G+\sigma)^2 \frac{E^2}{3}.
\end{aligned}
\tag{23}
\]

Since $\eta\le 1/L$, the coefficient of the last term is $O(\eta^3 L^2 E^2)$ and can be absorbed into the $\eta^2\sigma^2/(K E)$ term for a clean expression.  Dropping the non‑negative $G^2$ contribution (which only makes the bound looser) we obtain the simplified recursion
\[
\E_t\norm{w_{t+1}-w^\star}^2
\le \big(1+\eta\big)\norm{w_t-w^\star}^2
-2\eta\big(f(w_t)-f(w^\star)\big)
+ \eta\Delta^2
+ \frac{\eta^2\sigma^2}{K E}
+ 2\eta^2\Delta^2 .
\tag{24}
\]

\subsection*{Step 8: Telescoping over $T$ rounds}
Take total expectation and sum (24) from $t=0$ to $T-1$:
\[
\begin{aligned}
\E\norm{w_T-w^\star}^2
&\le (1+\eta)^T \norm{w_0-w^\star}^2
-2\eta\sum_{t=0}^{T-1}\E\big[f(w_t)-f(w^\star)\big] \\
&\quad + T\eta\Delta^2 + \frac{T\eta^2\sigma^2}{K E}+2T\eta^2\Delta^2 .
\end{aligned}
\tag{25}
\]

Discard the non‑negative left‑hand side and use $(1+\eta)^T\le e^{\eta T}\le 1+ \eta T$ for $\eta\le 1$:
\[
2\eta\sum_{t=0}^{T-1}\E\big[f(w_t)-f(w^\star)\big]
\le \norm{w_0-w^\star}^2 + T\eta\Delta^2 + \frac{T\eta^2\sigma^2}{K E}+2T\eta^2\Delta^2 .
\tag{26}
\]

Dividing by $2\eta T$ yields a bound on the average iterate $\bar w_T$:
\[
\begin{aligned}
\E\big[f(\bar w_T)-f(w^\star)\big]
&\le \frac{\norm{w_0-w^\star}^2}{2\eta K E T}
+ \frac{\eta L\sigma^2}{2K}
+ \frac{\eta L\Delta^2 (E-1)}{2},
\end{aligned}
\tag{27}
\]
where we used $\Delta^2$ together with the factor $(E-1)$ that appears when expanding the drift accumulation (Lemma~\ref{lem:drift}).  Equation (27) is exactly the statement (4) of Theorem~\ref{thm:conv}.

\subsection*{Step 9: Choice of stepsize}
Setting $\eta = \frac{c}{\sqrt{T}}$ (with $c\le 1/L$) in (27) gives
\[
\E\big[f(\bar w_T)-f(w^\star)\big]
= O\!\Big(\frac{1}{\sqrt{T}}+\frac{E-1}{\sqrt{T}}\Big),
\]
which completes the proof. \hfill $\square$

\section*{Rate Derivation (With Constants)}
From (27) we can write the explicit constant‑wise bound:
\[
\E\big[f(\bar w_T)-f(w^\star)\big]
\le
\underbrace{\frac{\norm{w_0-w^\star}^2}{2\eta K E T}}_{\text{optimization term}}
\;+\;
\underbrace{\frac{\eta L\sigma^2}{2K}}_{\text{stochastic variance}}
\;+\;
\underbrace{\frac{\eta L\Delta^2 (E-1)}{2}}_{\text{client‑drift penalty}} .
\]
Choosing $\eta = \min\big\{\frac{1}{L},\,\sqrt{\frac{\norm{w_0-w^\star}^2}{\sigma^2 K E T}}\big\}$ balances the first two terms and yields
\[
\E\big[f(\bar w_T)-f(w^\star)\big]
\le
\frac{L\norm{w_0-w^\star}^2}{2KET}
+ \frac{\Delta^2 (E-1)}{2L}.
\]

\section*{Special Cases}
\begin{enumerate}
\item \textbf{Homogeneous data ($\Delta=0$).}  The drift term vanishes and (27) reduces to the standard $O\!\big(\frac{1}{\eta KET}+\eta\sigma^2/K\big)$ bound.
\item \textbf{Single client ($K=1$).}  The variance term becomes $\frac{\eta L\sigma^2}{2}$, matching classical SGD.
\item \textbf{One local step ($E=1$).}  The drift penalty disappears; the algorithm is equivalent to parallel mini‑batch SGD.
\end{enumerate}

\end{document}